\documentclass[15pt]{ctexart}
\usepackage{xeCJK}
	\setCJKmainfont{SimSun}         % 缺省中文字体为宋体
	\setmainfont{Times New Roman}   % 缺省英文字体 Times New Roman
\usepackage{geometry}
	% \geometry{left=2.5cm,right=2.5cm,top=2.5cm,bottom=3.5cm}
	\geometry{bottom=3.5cm}
\usepackage{float}
\usepackage{listings}
\usepackage{graphicx}
\usepackage{appendix}
\usepackage[colorlinks]{hyperref}
\usepackage{amsmath}
\usepackage{indentfirst}
\usepackage{enumerate}
\usepackage{fancyhdr}
\usepackage{textcomp}
\usepackage{appendix} 
\usepackage{multirow}
\usepackage{geometry}
\geometry{verbose,letterpaper}
\usepackage{media9}
\pagestyle{fancy}
\lhead{虚拟路由}
\rhead{\thepage}
\cfoot{}
\lstset{frameshape={RYRYNYYYY}{yny}{yny}{RYRYNYYYY}, backgroundcolor=\color[RGB]{245,245,244}}

\begin{document}
\begin{titlepage}
    \centering
    \includegraphics[scale=0.9]{imgs/SYSULogo.png}\par\vspace{1cm}
    \vspace{1cm}
    {\scshape\huge 虚拟路由 \\ 项目报告 \\ \centering \scshape \Huge 模拟网络层路由 \par}
    \vspace{1.5cm}
    {\Large\bfseries \flushleft 学院：数据科学与计算机学院 \\ 专业：计算机科学与技术 \\ 年级：2016级 \\组长（学号）：王锡淮（16337236）\\组员（学号）：杨陈泽（16337271）\\
    组员（学号）：肖遥（16337258）\par}
    % \vspace{2cm}

% Bottom of the page
    % {\large \today\par}
\end{titlepage}
\tableofcontents
\newpage

\begin{table}[H]
	\centering
	\begin{tabular}{|c|}
		\hline
		Command \; \vline \; Version \; \vline \; Routing domain \; \\
		\hline
		Source address \\
		\hline
		Address family \; \vline \; Route tag \; \\
		\hline
		address \\
		\hline
		Netmask \\
		\hline
		Next hop \\
		\hline
		metric \\
		\hline
		repeat of last 21 bytes \\
		\hline
		$\cdots$ \\
		\hline 
	\end{tabular}	
\end{table}
报文项解释：
\begin{enumerate}
	\item Command（1 Byte）：指明这条报文的类型。
	\begin{enumerate}[]
		\item 0：普通报文。
		\item 1：请求报文。
		\item 2：响应（包含路由表）报文。
		% \item 3：广播request。
		% \item 4：广播response。
	\end{enumerate}
	\item Version（1 Byte）：RIP协议的版本。
	\item Routing domain（2 Bytes）：指明这是RIP中的哪一步。
	\item Source address（6 Bytes）：指明发送方的源地址和监听端口。
	\item Address family（2 Bytes）：使用什么作为地址，使用IP该项为2。
	\item Route tag（2 Bytes）：本项目中不用。
	\item address（6 Bytes）：ip和端口号。
	\item Netmask（4 Bytes）：子网掩码。
	\item Next hop（6 Bytes）：下一跳路由的ip和端口。
	\item matric（1 Bytes）：代价度量（跳数）。
\end{enumerate}
\begin{table}[H]
	\centering
	\begin{tabular}{|c|}
		\hline
		Command \; \vline \; Version \; \vline \; Routing domain \; \\
		\hline
		Source address \\
		\hline
		Address family \; \vline \; Route tag \; \\
		\hline
		address \\
		\hline
		Next hop \\
		\hline
		metric \\
		\hline
	\end{tabular}
\end{table}
报文项的内容和response表项一致。
\begin{table}[H]
	\centering
	\begin{tabular}{|c|}
		\hline
		Command \\
		\hline
		source address \\
		\hline
		destination address \\
		\hline
		payload \\
		\hline
	\end{tabular}		
\end{table}
其中的源地址和目的地址包含ip和监听端口号。
% \par 需要特别指明的是当用于广播时，目的ip写的是224.0.0.9。
\end{document}